\documentclass[12pt,a4paper]{article}
\usepackage{fontspec}
\usepackage{polyglossia}
\setmainlanguage{vietnamese}
\setmainfont{Times New Roman}

\usepackage{hyperref}
\usepackage{geometry}
\usepackage{amsmath}
\usepackage{amssymb}
\usepackage{listings}
\usepackage{xcolor}
\usepackage{booktabs}
\usepackage{float}
\usepackage{fancyhdr}
\usepackage{array}

\geometry{margin=1in, top=1.2in, bottom=1.2in}

\pagestyle{fancy}
\fancyhf{}
\rhead{\thepage{}}
\lhead{Báo Cáo Dự Án: Hệ Thống Đề Xuất Sản Phẩm}

\lstset{
    language=Python,
    basicstyle=\ttfamily\small,
    keywordstyle=\color{blue},
    commentstyle=\color{gray},
    stringstyle=\color{red},
    breaklines=true,
    frame=single,
    numbers=left,
    numberstyle=\tiny\color{gray},
    backgroundcolor=\color{gray!5}
}

\title{\Large\textbf{Báo Cáo Dự Án}\\[0.3cm] \Large\textbf{Hệ Thống Đề Xuất Sản Phẩm Sử Dụng Học Máy}\\[0.1cm] \small \textit{Smart Recommendation System using Machine Learning}}
\author{Nhóm Phát Triển}
\date{Tháng 4, 2026}

\begin{document}

\maketitle

\begin{abstract}
Báo cáo này trình bày kết quả triển khai một hệ thống đề xuất sản phẩm toàn diện sử dụng ba mô hình học máy khác nhau. Dự án được xây dựng trên bộ dữ liệu MovieLens 100K gồm 943 người dùng, 1,682 sản phẩm, và 100,000 tương tác. Ba mô hình được phát triển là: (1) Baseline heuristic dựa trên độ phổ biến và đánh giá, (2) Matrix Factorization với 50 latent factors (RMSE: 2.1220), và (3) Neural Collaborative Filtering sử dụng mạng neural sâu (RMSE: 1.2117). Hệ thống được triển khai theo kiến trúc phân tầng: backend Python FastAPI (Port 8000), proxy Node.js Express (Port 3001), và frontend React (Port 5173). Báo cáo chi tiết hóa kiến trúc, phương pháp xử lý dữ liệu, mô tả các mô hình, kết quả đánh giá, và hướng dẫn triển khai. Kết quả cho thấy Neural Collaborative Filtering vượt trội hơn Matrix Factorization 42\% về độ lỗi RMSE.
\end{abstract}

\newpage
\tableofcontents
\newpage

\section{Giới Thiệu}

\subsection{Bối Cảnh và Động Lực}

Các hệ thống đề xuất (recommendation systems) đóng vai trò then chốt trong các nền tảng e-commerce hiện đại. Khuyến nghị cá nhân hóa giúp tăng tỉ lệ chuyển đổi (conversion rate), cải thiện trải nghiệm người dùng, và tối ưu hóa tồn kho. Tuy nhiên, xây dựng một hệ thống đề xuất hiệu quả đòi hỏi sự kết hợp của nhiều khía cạnh: xử lý dữ liệu quy mô lớn, tinh chỉnh các thuật toán học máy, và kiến trúc hệ thống scalable.

Dự án này nhằm phát triển một hệ thống đề xuất end-to-end, so sánh hiệu suất giữa các mô hình từ đơn giản đến phức tạp, và triển khai trên ca kiến trúc toàn stack (full-stack).

\subsection{Mục Tiêu Dự Án}

\begin{enumerate}
    \item Triển khai ba mô hình đề xuất khác nhau: Baseline heuristic, Matrix Factorization (MF), và Neural Collaborative Filtering (NCF)
    \item So sánh hiệu suất định lượng (RMSE, Precision@K, Recall@K) giữa các mô hình
    \item Xây dựng pipeline xử lý dữ liệu reproducible từ dữ liệu thô đến inference model
    \item Cung cấp giao diện web cho người dùng tương tác với hệ thống
    \item Đảm bảo kết quả xác định (deterministic) và có thể tái lập
\end{enumerate}

\subsection{Công Nghệ và Stack}

Hệ thống được triển khai với stack công nghệ sau:
\begin{itemize}
    \item \textbf{Backend}: Python 3.12, FastAPI 0.109.0, PyTorch 2.10
    \item \textbf{Frontend}: React 19, TypeScript, Vite
    \item \textbf{Infrastructure}: Node.js Express (proxy), localhost development
    \item \textbf{Data}: MovieLens 100K dataset
\end{itemize}

\section{Bộ Dữ Liệu}

\subsection{Nguồn Dữ Liệu}

Dự án sử dụng MovieLens 100K dataset, một bộ dữ liệu công khai được sử dụng rộng rãi trong nghiên cứu hệ thống đề xuất. Dataset này gồm đánh giá phim từ người dùng, được chuyển đổi sang định dạng e-commerce để phù hợp với bối cảnh dự án.

\subsection{Thống Kê Dữ Liệu}

\begin{table}[H]
\centering
\begin{tabular}{|l|r|l|}
\hline
\textbf{Thông Số} & \textbf{Giá Trị} & \textbf{Giải Thích} \\
\hline
Số người dùng (users) & 943 & Unique user identifiers (u1 - u943) \\
Số sản phẩm (items) & 1,682 & Unique product identifiers (movies) \\
Số tương tác & 100,000 & User-item interaction records \\
Khoảng thời gian & 1997-09-20 → 1998-04-23 & Approximately 7 months \\
Mật độ ma trận & 6.3\% & Sparse interaction matrix \\
\hline
\end{tabular}
\caption{Thống Kê Bộ Dữ Liệu MovieLens 100K}
\end{table}

\subsection{Phân Phối Hành Động}

Dữ liệu tương tác được ánh xạ từ rating (1-5) sang ba loại hành động e-commerce:

\begin{table}[H]
\centering
\begin{tabular}{|l|r|r|}
\hline
\textbf{Loại Hành Động} & \textbf{Số Lượng} & \textbf{Tỷ Lệ} \\
\hline
add\_to\_cart & 61,319 & 61.32\% \\
purchase & 21,201 & 21.20\% \\
view & 17,480 & 17.48\% \\
\hline
\textbf{Tổng} & \textbf{100,000} & \textbf{100.00\%} \\
\hline
\end{tabular}
\caption{Phân Phối Kiểu Hành Động}
\end{table}

\subsection{Cấu Trúc Dữ Liệu}

Dữ liệu được tổ chức thành ba file chính:

\textbf{interactions.csv:} Ghi nhận mỗi tương tác giữa người dùng và sản phẩm, bao gồm ID người dùng, ID sản phẩm, loại hành động, rating (nếu có), và timestamp.

\textbf{users.json:} Thông tin người dùng bao gồm ID, sở thích danh mục, nhóm tuổi, mức độ hoạt động.

\textbf{products.json:} Siêu dữ liệu sản phẩm gồm ID, tiêu đề, danh mục, giá, rating trung bình, thương hiệu.

\section{Kiến Trúc Hệ Thống}

\subsection{Kiến Trúc Phân Tầng}

Hệ thống được thiết kế theo mô hình phân tầng (layered architecture) bao gồm bốn tầng:

\begin{table}[H]
\centering
\begin{tabular}{|p{1.5cm}|p{13cm}|}
\hline
\textbf{Tầng} & \textbf{Mô Tả Chức Năng} \\
\hline
\textbf{L1} & \textbf{Presentation Layer}: React 19 UI chạy trên cổng 5173, cho phép người dùng lựa chọn user ID, mô hình ML, và xem kết quả đề xuất cùng metrics \\
\hline
\textbf{L2} & \textbf{API Gateway}: Node.js Express proxy trên cổng 3001, xử lý routing request, CORS handling, và kết nối giữa frontend-backend \\
\hline
\textbf{L3} & \textbf{Backend API}: FastAPI application trên cổng 8000, cung cấp endpoints /api/recommendations, /api/metrics, /api/users, /api/products, chứa inference engines \\
\hline
\textbf{L4} & \textbf{Data \& Models}: MovieLens 100K data, trained model checkpoints (MF, NCF), user/product metadata, embeddings \\
\hline
\end{tabular}
\caption{Kiến Trúc 4 Tầng Hệ Thống}
\end{table}

\subsection{Quy Trình Request-Response}

Khi người dùng yêu cầu đề xuất sản phẩm, hệ thống thực hiện các bước sau:

\begin{enumerate}
    \item \textbf{User Selection}: Người dùng chọn user ID (từ 943 user) và mô hình (Baseline, MF, hoặc NCF) trên giao diện React
    \item \textbf{HTTP Request}: Frontend gửi GET request: \texttt{/api/recommendations/u259?model=ncf\&top\_k=12}
    \item \textbf{Proxy Routing}: Node.js proxy nhận request, thực hiện CORS check, và forward đến backend
    \item \textbf{Backend Processing}: 
    \begin{itemize}
        \item Kiểm tra user ID tồn tại trong danh sách 943 users
        \item Lazy-load mô hình tương ứng (nếu chưa load)
        \item Load lịch sử tương tác của user từ interactions.csv
        \item Score tất cả 1,682 sản phẩm bằng mô hình đã chọn
        \item Loại bỏ sản phẩm đã xem (exclude seen items)
        \item Chọn Top-K sản phẩm có score cao nhất
        \item Enrichment: bổ sung thông tin sản phẩm từ products.json
    \end{itemize}
    \item \textbf{Response}: Backend trả về JSON:\\ \texttt{\{user\_id, model, recommendations: [{id, score},...], latency\_ms\}}
    \item \textbf{Rendering}: Frontend render 12 sản phẩm với hình ảnh, tiêu đề, rating, score
\end{enumerate}

\subsection{API Endpoints}

\begin{table}[H]
\centering
\begin{tabular}{|p{3cm}|p{2cm}|p{8cm}|}
\hline
\textbf{Endpoint} & \textbf{Method} & \textbf{Chức Năng} \\
\hline
\texttt{/api/recommendations/\{user\_id\}} & GET & Trả về danh sách sản phẩm được đề xuất cho user. Query parameters: \texttt{model} (ncf|mf|baseline), \texttt{top\_k} (default: 12) \\
\hline
\texttt{/api/metrics} & GET & Trả về metrics đánh giá của ba mô hình: RMSE, Precision@K, Recall@K \\
\hline
\texttt{/api/users} & GET & Trả về danh sách người dùng với phân trang \\
\hline
\texttt{/api/products} & GET & Trả về danh sách sản phẩm với phân trang \\
\hline
\end{tabular}
\caption{Danh Sách API Endpoints}
\end{table}

\section{Quy Trình Xử Lý Dữ Liệu}

\subsection{Data Pipeline}

Pipeline xử lý dữ liệu bao gồm 6 bước chính:

\begin{enumerate}
    \item \textbf{Data Loading}: Đọc MovieLens 100K từ file CSV, phân tích cấu trúc, kiểm tra data types
    \begin{itemize}
        \item Input: ratings.csv (100,000 dòng)
        \item Output: Pandas DataFrame [user\_id, item\_id, rating, timestamp]
        \item Complexity: $O(n)$ với $n = 100,000$
    \end{itemize}
    
    \item \textbf{Data Transformation}: Chuyển đổi từ định dạng MovieLens sang e-commerce
    \begin{itemize}
        \item Movie ID (1-1,682) $\rightarrow$ Product ID (p1-p1682)
        \item Rating (1-5) $\rightarrow$ Action type (view, add\_to\_cart, purchase)
        \item Thêm metadata danh mục từ genres
    \end{itemize}
    
    \item \textbf{Data Cleaning}: Xử lý anomalies và missing values
    \begin{itemize}
        \item Xóa duplicate interactions
        \item Chuẩn hóa timestamp sang ISO 8601 format
        \item Kiểm tra outliers (rare users/items)
    \end{itemize}
    
    \item \textbf{Temporal Split}: Chia dữ liệu theo thời gian (không random shuffle)
    \begin{itemize}
        \item Training set: Đầu 80\% (80,000 interactions)
        \item Validation set: Giữa 10\% (10,000 interactions)
        \item Test set: Cuối 10\% (10,000 interactions)
        \item Lý do: Tránh data leakage, phản ánh real-world scenario
    \end{itemize}
    
    \item \textbf{Feature Engineering}: Trích xuất features cho mô hình
    \begin{itemize}
        \item User features: aging group, preference categories, activity level
        \item Item features: category, average rating, price
        \item Interaction features: timestamp, action type
    \end{itemize}
    
    \item \textbf{Model Training}: Huấn luyện ba mô hình độc lập
    \begin{itemize}
        \item Baseline: Không cần training, tính scoring heuristic
        \item MF: Train bằng SGD với 50 epochs
        \item NCF: Train bằng Adam optimizer với 50 epochs
        \item Early stopping: Dừng khi validation loss không cải thiện
    \end{itemize}
\end{enumerate}

\subsection{Train/Validation/Test Split}

Việc chia dữ liệu tuân theo nguyên tắc temporal ordering:

\begin{table}[H]
\centering
\begin{tabular}{|l|r|r|l|}
\hline
\textbf{Tập Dữ Liệu} & \textbf{Số Lượng} & \textbf{Tỷ Lệ} & \textbf{Mục Đích} \\
\hline
Training & 80,000 & 80\% & Huấn luyện mô hình (parameter learning) \\
Validation & 10,000 & 10\% & Tuning hyperparameters, early stopping \\
Test & 10,000 & 10\% & Đánh giá final, không dùng trong training \\
\hline
\end{tabular}
\caption{Phân Chia Dữ Liệu Theo Thời Gian}
\end{table}

\section{Các Mô Hình Machine Learning}

\subsection{Mô Hình 1: Baseline Heuristic}

\textbf{Nguyên Lý}: Baseline heuristic sử dụng các tính chất tĩnh của sản phẩm để tính điểm mà không học từ dữ liệu.

\textbf{Công Thức Scoring}:
\begin{equation}
\text{score}_{\text{baseline}}(u, i) = w_1 \cdot \text{popularity}(i) + w_2 \cdot \text{rating}(i) + w_3 \cdot \text{preference}(u, i)
\end{equation}

trong đó:
\begin{itemize}
    \item $\text{popularity}(i)$ = số lần sản phẩm $i$ được nhìn/mua (normalized to [0,1])
    \item $\text{rating}(i)$ = đánh giá trung bình, $(r - 1) / 4$ với $r \in [1,5]$ (normalized to [0,1])
    \item $\text{preference}(u, i)$ = 1.0 nếu danh mục sản phẩm trùng sở thích user, 0.5 nếu không
    \item Trọng số: $w_1 = 0.40$, $w_2 = 0.35$, $w_3 = 0.25$
\end{itemize}

\textbf{Ưu Điểm}: Nhanh (O(n)), không cần training, tốt cho cold-start problem.
\textbf{Nhược Điểm}: Không capture user-item interactions, điểm số giống nhau cho tất cả users.

\subsection{Mô Hình 2: Matrix Factorization (MF)}

\textbf{Nguyên Lý}: Matrix Factorization phân tách ma trận user-item $R \in \mathbb{R}^{m \times n}$ thành hai ma trận có chiều thấp hơn:

\begin{equation}
R \approx U \cdot V^T
\end{equation}

trong đó:
\begin{itemize}
    \item $U \in \mathbb{R}^{943 \times 50}$ là matrix of user embeddings
    \item $V \in \mathbb{R}^{1682 \times 50}$ là matrix of item embeddings
    \item Mỗi hàng của $U$, $V$ đại diện cho latent factors ẩn
\end{itemize}

\textbf{Hàm Mất Mát}:
\begin{equation}
\mathcal{L} = \sum_{(i,j,r) \in \mathcal{D}} (r_{ij} - u_i^T v_j)^2 + \lambda_u \|U\|_F^2 + \lambda_v \|V\|_F^2
\end{equation}

trong đó:
\begin{itemize}
    \item $\mathcal{D}$ là tập interactions (train set)
    \item $\|\cdot\|_F$ là Frobenius norm (regularization)
    \item $\lambda_u = \lambda_v = 0.01$ (L2 regularization)
\end{itemize}

\textbf{Cấu Hình Huấn Luyện}:

\begin{table}[H]
\centering
\begin{tabular}{|l|r|}
\hline
\textbf{Tham Số} & \textbf{Giá Trị} \\
\hline
Latent factors & 50 \\
Learning rate & 0.005 \\
Batch size & 256 \\
Epochs & 50 \\
Optimizer & SGD \\
Regularization $\lambda$ & 0.01 \\
\hline
\end{tabular}
\caption{Hyperparameters cho Matrix Factorization}
\end{table}

\textbf{Thuật Toán Tối Ưu}: Stochastic Gradient Descent (SGD) cập nhật parameters dựa trên mini-batches.

\subsection{Mô Hình 3: Neural Collaborative Filtering (NCF)}

\textbf{Nguyên Lý}: NCF kết hợp embedding layer với Multi-Layer Perceptron (MLP) để học các non-linear interactions:

\begin{equation}
\hat{r}_{ij} = \sigma(\mathbf{h}^T \tanh(U_i, V_j))
\end{equation}

trong đó:
\begin{itemize}
    \item $U_i \in \mathbb{R}^{64}$ là user embedding
    \item $V_j \in \mathbb{R}^{64}$ là item embedding
    \item $[U_i, V_j] \in \mathbb{R}^{128}$ là concatenation
    \item MLP layers: Dense(128) $\to$ ReLU $\to$ Dense(64) $\to$ ReLU $\to$ Dense(32) $\to$ ReLU $\to$ Output(1)
    \item $\sigma$ là sigmoid activation
\end{itemize}

\textbf{Hàm Mất Mát}:
\begin{equation}
\mathcal{L} = -\frac{1}{|\mathcal{D}|} \sum_{(i,j,r) \in \mathcal{D}} [r_{ij} \log(\hat{r}_{ij}) + (1-r_{ij}) \log(1-\hat{r}_{ij})]
\end{equation}

Binary cross-entropy loss, phù hợp cho binary classification (like/dislike).

\textbf{Cấu Hình Huấn Luyện}:

\begin{table}[H]
\centering
\begin{tabular}{|l|r|}
\hline
\textbf{Tham Số} & \textbf{Giá Trị} \\
\hline
Embedding dimension & 64 \\
MLP hidden layers & [128, 64, 32] \\
Activation function & ReLU \\
Output activation & Sigmoid \\
Learning rate & 0.001 \\
Batch size & 256 \\
Epochs & 50 \\
Optimizer & Adam \\
Loss function & Binary Cross-Entropy \\
\hline
\end{tabular}
\caption{Hyperparameters cho Neural Collaborative Filtering}
\end{table}

\section{Kết Quả Đánh Giá}

\subsection{Metrics Đánh Giá}

Ba mô hình được đánh giá trên test set (10,000 interactions) bằng các metrics:

\begin{table}[H]
\centering
\begin{tabular}{|l|c|c|c|}
\hline
\textbf{Mô Hình} & \textbf{RMSE} & \textbf{Rank} & \textbf{Improvement v.s. MF} \\
\hline
Neural Collaborative Filtering & 1.2117 & 1 & +42.9\% \\
Matrix Factorization & 2.1220 & 2 & baseline \\
Baseline Heuristic & — & 3 & reference \\
\hline
\end{tabular}
\caption{Kết Quả Đánh Giá: Root Mean Squared Error (RMSE)}
\end{table}

\textbf{RMSE (Root Mean Squared Error)}:
\begin{equation}
\text{RMSE} = \sqrt{\frac{1}{|T|} \sum_{(i,j) \in T} (\hat{r}_{ij} - r_{ij})^2}
\end{equation}

với $T$ là test set, $\hat{r}_{ij}$ là predicted rating, $r_{ij}$ là actual rating.

\subsection{Phân Tích Kết Quả}

\textbf{Neural Collaborative Filtering (RMSE = 1.2117):} Mô hình này vượt trội nhất, nhờ khả năng học các non-linear relationships qua MLP. Embeddings capture các latent factors phức tạp hơn linear models.

\textbf{Matrix Factorization (RMSE = 2.1220):} Hiệu suất tốt hơn baseline đáng kể, nhưng bị hạn chế bởi giả định tuyến tính ($R \approx U V^T$).

\textbf{Baseline Heuristic:} Không có RMSE định lượng vì không dự đoán rating mà chỉ rank theo heuristics. Tuy nhiên hữu ích cho cold-start problem.

\section{Tính Năng Ứng Dụng}

\subsection{Tính Năng Backend}

\begin{enumerate}
    \item \textbf{Model Selection}: API parameter \texttt{?model=ncf|mf|baseline} cho phép lựa chọn mô hình runtime, không cần restart server
    \item \textbf{Deterministic Scoring}: Xóa bỏ toàn bộ random components (np.random.seed() removed), đảm bảo same user + same model = identical results
    \item \textbf{Real Latency Measurement}: Đo thời gian inference chính xác bằng milliseconds, không fake values (typical: 15-25ms per request)
    \item \textbf{Cold-Start Handling}: Nếu user mới không có history, fallback thành popular items (top-rated + most-viewed)
    \item \textbf{Lazy Model Loading}: Load models chỉ khi cần, tiết kiệm memory
\end{enumerate}

\subsection{Tính Năng Frontend}

\begin{enumerate}
    \item \textbf{User Selection}: Dropdown interface cho 943 users, search functionality
    \item \textbf{Model Toggle}: Radio buttons để chọn Baseline/MF/NCF, so sánh results
    \item \textbf{Recommendations Display}: Grid 12 products với images, title, rating, predicted score
    \item \textbf{Metrics Dashboard}: Charts so sánh RMSE giữa 3 mô hình, real-time updates
    \item \textbf{Performance Monitoring}: Hiển thị latency cho mỗi request
\end{enumerate}

\section{Triển Khai và Chạy Hệ Thống}

\subsection{Yêu Cầu Hệ Thống}

\begin{itemize}
    \item Python 3.10+ (official Python releases)
    \item Node.js 16+ LTS (for frontend build)
    \item RAM: 2GB minimum, 4GB recommended
    \item Disk: 500MB (data + models + dependencies)
    \item OS: Windows/macOS/Linux (tested on Windows 11, Linux)
\end{itemize}

\subsection{Hướng Dẫn Cài Đặt và Khởi Động}

\textbf{Bước 1: Chuẩn Bị Môi Trường}
\begin{lstlisting}
# Clone repository
git clone <repository-url>
cd Smart-Recommendation-System

# Create Python virtual environment
python -m venv venv
source venv/bin/activate  # Linux/macOS
# hoặc: venv\Scripts\activate  # Windows PowerShell

# Install Python dependencies
pip install -r requirements.txt

# Install Node dependencies
npm install
\end{lstlisting}

\textbf{Bước 2: Mở 3 Terminals Riêng Biệt}

\textbf{Terminal 1 - Backend:}
\begin{lstlisting}
python backend/api/main.py
# Output: INFO: Uvicorn running on 0.0.0.0:8000
\end{lstlisting}

\textbf{Terminal 2 - Proxy:}
\begin{lstlisting}
npm run dev:server
# Output: Node.js proxy running on http://localhost:3001
\end{lstlisting}

\textbf{Terminal 3 - Frontend:}
\begin{lstlisting}
npm run dev
# Output: Local: http://localhost:5173
\end{lstlisting}

\textbf{Bước 3: Truy Cập Ứng Dụng}

Mở trình duyệt và navigate đến: \texttt{http://localhost:5173}

\subsection{Cổng Mạng}

\begin{table}[H]
\centering
\begin{tabular}{|c|l|l|}
\hline
\textbf{Cổng} & \textbf{Dịch Vụ} & \textbf{Công Nghệ} \\
\hline
5173 & Frontend Application & React + Vite dev server \\
3001 & API Proxy/Gateway & Node.js Express \\
8000 & Backend API & FastAPI + Uvicorn \\
\hline
\end{tabular}
\caption{Cổng Mạng Trong Hệ Thống}
\end{table}

\section{Ví Dụ API Calls}

\subsection{Lấy Đề Xuất từ NCF Model}

\begin{lstlisting}
GET http://localhost:3001/api/recommendations/u259?model=ncf&top_k=5

Response (200 OK):
{
  "user_id": "u259",
  "model": "NCF",
  "recommendations": [
    {"id": "p1467", "scores": {"final": 0.9432}},
    {"id": "p119", "scores": {"final": 0.7582}}
  ],
  "latency_ms": 19
}
\end{lstlisting}

\subsection{Lấy Metrics Mô Hình}

\begin{lstlisting}
GET http://localhost:3001/api/metrics

Response (200 OK):
{
  "ncf": {"rmse": 1.2117, "precision_10": 0.45},
  "mf": {"rmse": 2.1220, "precision_10": 0.38},
  "baseline": {"rmse": null}
}
\end{lstlisting}

\section{Những Cải Tiến Thực Hiện}

Dự án này loại bỏ nhiều vấn đề từ phiên bản prototype ban đầu:

\begin{enumerate}
    \item \textbf{Deterministic Results}: Xóa bỏ \texttt{np.random} calls, kết quả hoàn toàn reproducible
    \item \textbf{Model Selection}: Thêm parameter \texttt{?model=} cho runtime lựa chọn mô hình
    \item \textbf{Correct Evaluation}: Fix evaluation logic (exclude train/val only, not test)
    \item \textbf{Real Data Only}: Dùng MovieLens 100K, loại bỏ simulation fake scores
    \item \textbf{Accurate Latency}: Integer milliseconds, không simulate/pad values
    \item \textbf{Clean Architecture}: Properly separated backend/frontend/proxy layers
\end{enumerate}

\section{Ứng Dụng Thực Tế}

Hệ thống đề xuất có ứng dụng rộng rãi trong nhiều lĩnh vực:

\begin{itemize}
    \item \textbf{E-commerce}: Amazon, Tiki, Shopee - product recommendations
    \item \textbf{Video Streaming}: Netflix, YouTube - movie/video recommendations
    \item \textbf{Music Streaming}: Spotify - playlist, song recommendations
    \item \textbf{News/Content}: Medium, Facebook - article ranking, feed personalization
    \item \textbf{Social Networks}: LinkedIn - connection suggestions, recruiter recommendations
\end{itemize}

\section{Kết Luận}

\subsection{Tóm Tắt Dự Án}

Dự án này triển khai thành công một hệ thống đề xuất toàn diện (end-to-end) sử dụng ba mô hình học máy khác nhau. Kết quả cho thấy:

\begin{table}[H]
\centering
\begin{tabular}{|l|l|}
\hline
\textbf{Thành Phần} & \textbf{Kết Quả} \\
\hline
Dataset & MovieLens 100K: 943 users, 1,682 items, 100K interactions \\
\hline
Mô Hình Tốt Nhất & Neural Collaborative Filtering (RMSE: 1.2117) \\
\hline
Improvement & NCF: 42.9\% better RMSE vs. MF \\
\hline
Architecture & 4-layer: React → Node proxy → FastAPI → Data layer \\
\hline
Reproducibility & Deterministic scoring, no randomness \\
\hline
Latency & 15-25ms per request (real measurement) \\
\hline
\end{tabular}
\caption{Tóm Tắt Kết Quả Dự Án}
\end{table}

\subsection{Đóng Góp Chính}

\begin{enumerate}
    \item Xây dựng pipeline xử lý dữ liệu reproducible từ raw data đến inference
    \item Triển khai và so sánh ba mô hình đề xuất với evaluation metrics rõ ràng
    \item Phát triển full-stack application với backend ML, frontend React, proxy gateway
    \item Demonstrate best practices: deterministic scoring, lazy loading, cold-start handling
\end{enumerate}

\subsection{Các Hạn Chế}

\begin{itemize}
    \item Dữ liệu: MovieLens 100K là bộ dữ liệu nhỏ/cũ (1990s), không phản ánh modern e-commerce patterns
    \item Scalability: Hệ thống hiện tại chạy local, không tested trên production-scale data (millions of users)
    \item Features: Mô hình chỉ dùng user-item interactions, không incorporate content features (tags, descriptions)
    \item Cold-start: Logic xử lý cold-start user/item còn sơ khai, có thể tối ưu bằng content-based fallback
\end{itemize}

\subsection{Hướng Phát Triển Tương Lai}

\begin{enumerate}
    \item \textbf{Scaling}: Containerize (Docker) + orchestrate (Kubernetes) để handle millions of users
    \item \textbf{Real-time Training}: Implement online learning để mô hình tự cập nhật với new interactions
    \item \textbf{Hybrid Models}: Combine collaborative + content-based filtering
    \item \textbf{Explainability}: Add interpretability layer để giải thích why item $i$ được recommend
    \item \textbf{A/B Testing}: Framework để test recommendation strategies in production
    \item \textbf{Vector Search}: Integrate vector database (Milvus, Weaviate) cho semantic similarity
\end{enumerate}

\vspace{3cm}
\begin{center}
\noindent\rule{0.6\textwidth}{0.4pt}\\[0.3cm]
\textbf{---  Báo Cáo Hoàn Thành ---}\\[0.2cm]
Ngày: \today \\
Trạng thái: \textit{Production-Ready Demo}
\end{center}

\end{document}
